% Part: computability
% Chapter: recursive-functions
% Section: non-pr-functions

\documentclass[../../../include/open-logic-section]{subfiles}

\begin{document}

\olfileid{cmp}{rec}{npr}
\olsection{دوال غير عودية بدائية}

ليست كل دالة يتأتى حسابها بحسب
الحدس عودية بدائية. ولبيان ذلك نتأمل إمكان سرد جميع الدوال
العودية البدائية الأحادية الحجة في قائمة ${\OLId{latin-ordinary}{f}}_{\OLMathNumeral{0}}$ و${\OLId{latin-ordinary}{f}}_{\OLMathNumeral{1}}$ و${\OLId{latin-ordinary}{f}}_{\OLMathNumeral{2}}$ و…،
يُحسب بها ${\OLId{latin-ordinary}{f}}_{\OLId{latin-ordinary}{x}}$ عند الدخل ${\OLId{latin-ordinary}{y}}$ بطريقة فعالة. ومعنى هذا
أن الدالة ${\OLId{latin-ordinary}{g}}({\OLId{latin-ordinary}{x}},{\OLId{latin-ordinary}{y}})$ التي تعريفها
\[
{\OLId{latin-ordinary}{g}}({\OLId{latin-ordinary}{x}},{\OLId{latin-ordinary}{y}}) = {\OLId{latin-ordinary}{f}}_{\OLId{latin-ordinary}{x}}({\OLId{latin-ordinary}{y}})
\]
قابلة للحساب. وحينئذ تكون الدالة الآتية قابلة للحساب كذلك:
\begin{eqnarray*}
{\OLId{latin-ordinary}{h}}({\OLId{latin-ordinary}{x}}) & = & {\OLId{latin-ordinary}{g}}({\OLId{latin-ordinary}{x}},{\OLId{latin-ordinary}{x}}) + {\OLMathNumeral{1}} \\
& = & {\OLId{latin-ordinary}{f}}_{\OLId{latin-ordinary}{x}}({\OLId{latin-ordinary}{x}}) +{\OLMathNumeral{1}}.
\end{eqnarray*}
ومهما أخذنا دالة عودية بدائية ${\OLId{latin-ordinary}{f}}_{\OLId{latin-ordinary}{i}}$، خالفتها ${\OLId{latin-ordinary}{h}}$؛ إذ تختلف قيمة ${\OLId{latin-ordinary}{f}}_{\OLId{latin-ordinary}{i}}$ عند ${\OLId{latin-ordinary}{i}}$ عن القيمة التي أعطتها الدالة القطرية. لذلك
تكون ${\OLId{latin-ordinary}{h}}$ قابلة للحساب، ولا تكون عودية بدائية. وكذلك الحكم
في ${\OLId{latin-ordinary}{g}}$. فهذا إجراء «فعال» لحجة كانتور القطرية.

ولنا أمثلة أصرح على اجتماع قابلية الحساب وانتفاء العودية البدائية. فليكن
الرمز ${\OLId{latin-ordinary}{g}}^{\OLId{latin-ordinary}{n}}({\OLId{latin-ordinary}{x}})$، مثلًا، على ${\OLId{latin-ordinary}{g}}({\OLId{latin-ordinary}{g}}(\dots {\OLId{latin-ordinary}{g}}({\OLId{latin-ordinary}{x}})))$، وفيه ${\OLId{latin-ordinary}{n}}$ من ورود
${\OLId{latin-ordinary}{g}}$؛ ولنعرف متتالية الدوال ${\OLId{latin-ordinary}{g}}_{\OLMathNumeral{0}},{\OLId{latin-ordinary}{g}}_{\OLMathNumeral{1}},\dots$ بـ
\begin{eqnarray*}
{\OLId{latin-ordinary}{g}}_{\OLMathNumeral{0}}({\OLId{latin-ordinary}{x}}) & = & {\OLId{latin-ordinary}{x}}+{\OLMathNumeral{1}} \\
{\OLId{latin-ordinary}{g}}_{{\OLId{latin-ordinary}{n}} + {\OLMathNumeral{1}}}({\OLId{latin-ordinary}{x}}) & = & {\OLId{latin-ordinary}{g}}_{\OLId{latin-ordinary}{n}}^{\OLId{latin-ordinary}{x}}({\OLId{latin-ordinary}{x}})
\end{eqnarray*}
تحقّق أن كل ${\OLId{latin-ordinary}{g}}_{\OLId{latin-ordinary}{n}}$ عودية بدائية. ثم لاحظ أن نمو كل دالة تالية
يفوق نمو سابقتها كثيرًا؛ فـ${\OLId{latin-ordinary}{g}}_{\OLMathNumeral{1}}({\OLId{latin-ordinary}{x}})$ تساوي ${\OLMathNumeral{2}}{\OLId{latin-ordinary}{x}}$، و${\OLId{latin-ordinary}{g}}_{\OLMathNumeral{2}}({\OLId{latin-ordinary}{x}})$ تساوي
${\OLMathNumeral{2}}^{\OLId{latin-ordinary}{x}}\cdot {\OLId{latin-ordinary}{x}}$، وتنمو ${\OLId{latin-ordinary}{g}}_{\OLMathNumeral{3}}({\OLId{latin-ordinary}{x}})$ نموًا يقارب برجًا أسيًا يتألف من ${\OLId{latin-ordinary}{x}}$
نسخ من العدد~${\OLMathNumeral{2}}$. ودالة أكرمان--بيتر هي في جوهرها الدالة ${\OLId{latin-looped}{G}}({\OLId{latin-ordinary}{x}})={\OLId{latin-ordinary}{g}}_{\OLId{latin-ordinary}{x}}({\OLId{latin-ordinary}{x}})$،
ويثبت أن نموها يجاوز نمو كل دالة عودية بدائية.

ونرجع الآن إلى تعداد الدوال العودية البدائية. فقد جعلنا
لكل دالة منها ترميزًا رمزيًا؛ فيكفينا أن نعدد هذه الترميزات.
وإسناد عدد طبيعي $\#({\OLId{latin-looped}{F}})$ إلى كل ترميز ${\OLId{latin-looped}{F}}$ يتم عوديًا كما يأتي:
\begin{eqnarray*}
\#({\OLMathNumeral{0}}) & = & \langle {\OLMathNumeral{0}} \rangle \\
\#({\OLId{latin-looped}{S}}) & = & \langle {\OLMathNumeral{1}} \rangle \\
\#(\Proj{{\OLId{latin-ordinary}{n}}}{{\OLId{latin-ordinary}{i}}}) & = & \langle {\OLMathNumeral{2}}, {\OLId{latin-ordinary}{n}}, {\OLId{latin-ordinary}{i}} \rangle \\
\#(\fn{Comp}_{{\OLId{latin-ordinary}{k}},{\OLId{latin-ordinary}{l}}}[{\OLId{latin-looped}{H}},{\OLId{latin-looped}{G}}_{\OLMathNumeral{0}},\dots,{\OLId{latin-looped}{G}}_{{\OLId{latin-ordinary}{k}}-{\OLMathNumeral{1}}}]) & = & \langle
{\OLMathNumeral{3}},{\OLId{latin-ordinary}{k}},{\OLId{latin-ordinary}{l}},\#({\OLId{latin-looped}{H}}),\#({\OLId{latin-looped}{G}}_{\OLMathNumeral{0}}),\dots,\#({\OLId{latin-looped}{G}}_{{\OLId{latin-ordinary}{k}}-{\OLMathNumeral{1}}}) \rangle \\
\#(\fn{Rec}_{\OLId{latin-ordinary}{l}}[{\OLId{latin-looped}{G}},{\OLId{latin-looped}{H}}]) & = & \langle {\OLMathNumeral{4}}, {\OLId{latin-ordinary}{l}}, \#({\OLId{latin-looped}{G}}), \#({\OLId{latin-looped}{H}}) \rangle
\end{eqnarray*}
وقد استعملنا أن متتالية الأعداد تؤخذ عددًا
طبيعيًا، بحسب الترميز المتقدم. فصار لكل ترميز رمزي
عدد طبيعي. وليس كل متتالية، ولا كل عدد، رمزًا لترميز
صحيح؛ غير أن هذا لا يفسد التعداد، إذ نجعل ${\OLId{latin-ordinary}{f}}_{\OLId{latin-ordinary}{i}}$ الدالة العودية البدائية الأحادية
الحجة ذات الترميز الذي رمزه ${\OLId{latin-ordinary}{i}}$ إذا كان ${\OLId{latin-ordinary}{i}}$ يرمز ترميزًا كهذا، ونجعلها
الدالة الثابتة ${\OLMathNumeral{0}}$ في غير ذلك. فبهذا حصلت طريقة صريحة
لتعداد الدوال العودية البدائية الأحادية الحجة.

(وقد تتكرر الدالة الواحدة في القائمة، كالدالة الثابتة صفرًا. وليس
التكرار ناشئًا من اختيارنا رموز الأعداد وحده؛ فإن للدالة الثابتة صفرًا
نفسها أكثر من ترميز رمزي. وسيأتي أن دفع هذه
التكرارات لا يتأتى بطريقة قابلة للحساب؛ فمن ذلك أنه لا توجد دالة قابلة للحساب تحكم
هل يمثل ترميز معطى الدالة الثابتة صفرًا أم لا.)

فلنأخذ الآن ${\OLId{latin-ordinary}{g}}({\OLId{latin-ordinary}{x}},{\OLId{latin-ordinary}{y}})$ معرّفة بـ${\OLId{latin-ordinary}{f}}_{\OLId{latin-ordinary}{x}}({\OLId{latin-ordinary}{y}})$، ولتكن ${\OLId{latin-ordinary}{f}}_{\OLId{latin-ordinary}{x}}$ من
التعداد المتقدم. وبأي وجه نعلم قابلية ${\OLId{latin-ordinary}{g}}({\OLId{latin-ordinary}{x}},{\OLId{latin-ordinary}{y}})$ للحساب؟ وجهه بحسب
الحدس أننا لحساب ${\OLId{latin-ordinary}{g}}({\OLId{latin-ordinary}{x}},{\OLId{latin-ordinary}{y}})$ نحلل الرمز ${\OLId{latin-ordinary}{x}}$، ونتبين هل يدل على دالة
أحادية الحجة؛ فإن دل عليها حسبنا قيمتها عند الدخل~${\OLId{latin-ordinary}{y}}$.

\begin{tagblock}{TMs}
\begin{digress}
ولعلك ترى إمكان كتابة برنامج يؤدي ذلك، وإن احتاج إلى شيء من العمل،
بلغة جافا أو ++C مثلًا. فلنا حينئذ أن نستند إلى أطروحة
تشرش--تورنغ: كل ما يقبل الحساب بحسب الحدس يتأتى
حسابه بآلة تورنغ.

وأقرب إلى البرهان المباشر على قابلية ${\OLId{latin-ordinary}{g}}({\OLId{latin-ordinary}{x}},{\OLId{latin-ordinary}{y}})$ للحساب أن نصف آلة
تورنغ تحسبها وصفًا صريحًا؛ فنستغني عن أطروحة تشرش--تورنغ
وعن الرجوع إلى الحدس. وسنهيئ قريبًا ما يكفي من الأدوات لإثبات
قابلية ${\OLId{latin-ordinary}{g}}({\OLId{latin-ordinary}{x}},{\OLId{latin-ordinary}{y}})$ للحساب بواسطة نموذج يمكن
\emph{محاكاته} على آلة تورنغ، أعني الدوال العودية.
\end{digress}
\end{tagblock}

\end{document}
